\section{What the Full Framework Adds}\label{sec:framework}

The tour ran eight products through one machine and watched it hold. A book that
runs this machine for real needs more than the machine, and it would be dishonest
to close without saying what. This section names that additional machinery. The
point worth making first is that none of it is a second architecture bolted to the
side. Each piece is a consequence of the door, the log, and the views already
built: a reporting obligation is a view, an industry mapping is a function over the
log, a correctness claim is a check that runs against the log, the timer service is
named infrastructure outside the trust boundary, and the standing scope limits are
the same boundary the ledger reconciles against rather than governs. Where this
account does not carry a piece to its full depth, it stops and says so, rather than
implying the depth is not there.

\subsection{Regulatory reporting is a projection, not a second record}

A trading book must report its activity --- trades, transactions, positions ---
to authorities that keep their own records. It is tempting to build this as a
parallel pipeline: a second store of trade data, shaped for the report, kept in
step with the ledger by a process that copies and reconciles. That design
reintroduces exactly the failure the whole system exists to remove, because now
there are two records of one activity and a way for them to disagree.

There is no need for it. The economic content of any such report --- what happened,
in what quantity, to what unit, when --- is already in the move stream, because
that stream is the record of what happened. A report is therefore one more view:
a fold that reads the committed log and emits the report's rows in the report's
shape. It computes; it never mutates state, never writes back, and holds no
authority the log does not already hold. If two reports of the same activity must
agree, they agree for the same reason the balance sheet and the profit-and-loss
statement agree --- they are two folds of one log, and there is only one log for
them to fold. What a report needs beyond the economic content is reference data:
the identifiers and classifications an authority requires --- a trade's assigned
identifier, a party's registered legal identifier, a counterparty's regulatory
class. Those are not facts the ledger forms; they are supplied at the boundary, by
the authorities that own them, and joined to the projection as it is built. The
ledger records the activity and projects it; it does not become the registrar of
the identifiers the report is filed under.

\subsection{Alignment to a common domain model}

The industry describes a trade, and the events that happen to it over its life, in
a shared domain model --- a common vocabulary of products and lifecycle events that
lets independent systems talk about the same trade and mean the same thing. The
architecture is built to align with that vocabulary rather than to compete with it.
The alignment is a structural correspondence, and it is concrete: the ledger's unit
corresponds to the model's product --- and, for a bespoke instrument, to the trade
itself --- and the map from the model's business events to the ledger's transactions
is required to be total and deterministic. A trade becoming a position, a coupon
falling due, a corporate action landing on a holding --- each is an event in the
model and a transaction at the door, and the map carries one to the other. Because
the map is a function of the event alone, any conforming map inherits what the
architecture already guarantees: it composes (a sequence of events maps to the sequence of transactions),
it conserves (the moves it emits net to zero like any others), it keeps move order,
and it is idempotent under the same keys the door already deduplicates. A corporate
action is the one place the map does more than translate: the event supplies both
halves of the transaction the ledger needs --- the entitlement moves and the
status write that advances the basis, with its declared conversion --- so that the
basis change, which a plainer mapping would discard as it had no field for it, is
carried as the first-class fact it is. This account states the alignment and its
shape; it does not carry the field-by-field mapping, and does not pretend to.

\subsection{The correctness claims are executable, not asserted}

Throughout, this design has made claims of correctness: that quantity is conserved,
that a price and a quantity meet in one basis and never two, that every product has
exactly one declared reaction to every corporate-action class, that a producer's
economics can be recomputed from the log and compared, that a dated obligation
reaches a terminal state by its deadline. A claim like that is worth only as much as
the reader's ability to check it, and the framework's discipline is that each such
claim is not a sentence in prose but a property a machine re-checks against the
reference behaviour. Every claim is a universally quantified statement --- for every
transaction, for every unit, for every boundary --- and each is carried as an
executable oracle: it is fed cases drawn from the closed enumerations of products and
events the design admits, and it either holds on every case or names the case it
fails on. Structural claims hold by construction, so their oracles confirm that the
types admit no counterexample; economic claims are checked, so their oracles are the
recomputation the witness performs. This is the concrete content, turned on the
framework's own claims, of the discipline that a producer's output is trusted for
nothing and verifiable in everything: a claim in this framework is something a
machine re-checks, not something a reader is asked to take on trust, and coverage of
the event vocabulary is a finite checklist rather than a hope that testing found the
corner.

\subsection{The execution engine and the liveness it underwrites}

A lifecycle contract does not wake itself. Something must hold the dated
obligations of every open position --- an expiry, a coupon, an exercise window ---
and fire each when its time comes. That something is an execution engine, and it is
infrastructure outside the trust boundary, named apart from the door, the log, and
the views, holding no privileged access of any kind. It is bound by named
requirements and nothing more: it must fire durably, so a timer survives a restart;
it must deliver at least once, so a fire is never simply lost; and its firing is
delivery, not data --- it says \emph{when}, while everything the event means is read
from the log. Because the trigger is delivery and never data, an engine that stalls,
fails over, fires early, or fires twice can delay a discharge but cannot corrupt the
log: a late delivery submits the identical transaction, a duplicate delivery is
absorbed once at the door, and an early or spurious fire, finding its observation or
due-state not yet in the log, yields a refusal rather than a commit.

This is what lets the design carry liveness as a guarantee rather than a hope. Every
contractual obligation is a first-class logged object, registered atomically with
the position it binds and carrying three things: a deadline, a discharge predicate
on ledger state, and a compensation action. The guarantee is that by its deadline
every obligation reaches one of three terminal states --- Discharged, Compensated,
or Defaulted. The transition is itself an ordinary status write through the single
writer for obligation state: the discharge predicate met on the discharging
transaction sets Discharged; and once the deadline passes unmet, the engine wakes the
obligation's compensation contract, which sets Compensated when the compensation
action completes and Defaulted when even that cannot. No other component can move an
obligation to a terminal state, because obligation state, like every other fact, has
one legal writer. A coupon that should fire and does not is then not silence but a
visible object past its deadline, which any fold of the log can see; a dated duty
cannot be quietly forgotten, because forgetting it would be a state, and the state
is detectable. The honest content of the guarantee is where it is conditional. It
holds provided the engine is available at the deadline to deliver the trigger, the
executor that runs the woken contract is available to act, any external cooperation
the discharge requires is forthcoming, and the lifecycle function computes the right
discharge from what it reads. Those conditions are named because they are real: the
ledger cannot make a counterparty pay, and cannot fire a timer while the engine that
holds it is down. Under an engine outage the obligation is reached \emph{eventually},
when delivery resumes --- not instantaneously. What the log guarantees unconditionally
is that it never loses the obligation; it only waits for the fire, and the wait is
itself a recorded, detectable state.

\subsection{What this framework does not do}

Honesty about depth requires honesty about the edges. The framework has standing
limits, and they are not gaps to be filled later but boundaries it is built to
reconcile against rather than to govern. It does not form prices: what a structured
note is worth is decided by a model on the far side of a wall the ledger never
crosses, and the ledger records that number as a stamped observation without
reproducing or validating it. It is not the settlement infrastructure: it projects
committed activity outward as instructions, and the accounts that move real cash and
securities are elsewhere. It is not a legal-agreement system: a contract in this
design is the machine-readable form of a confirmation, not the confirmation's legal
force. It is not a regulatory-submission channel: it computes the reporting view, and
the filing of it is an external act. It is not a reference-data authority: the
identifiers, classifications, and registrations it consumes are owned by others.
Against each of these it reconciles at its boundary --- the external record is
independent, and reconciliation against it stays necessary, made detectable and
narrow rather than eliminated.

Two residual trust assumptions remain, and naming them is part of the honesty. The
first is at the ingest stamp: when market data enters, the stamp that says which
basis the datum was disseminated in is trusted to be true of the source, and the
architecture bounds what that trust buys --- a datum crosses a boundary only by a
declared operator, and a registered kind carries a mandatory invariance witness that
pins exactly what the kind's declaration is trusted for, no wider. The second is the
one the two-layer discipline has already been plain about: a contract's output is
admitted before its economics are recomputed. A transaction that is well-formed but
economically wrong --- pays \EUR{98} where \EUR{100} was owed --- is admitted, and
its error is found afterward by the witness and repaired by a compensating
transaction that names what it supersedes. Detection is therefore eventual, not
preventive. What is guaranteed absolutely is that the structural invariants never
break and that the wrong number is always findable from the log and detected whenever
the witness runs; what is not claimed is that the wrong number is stopped before it is
booked. How promptly the witness runs is a scheduling requirement the deployment owns,
standing alongside the engine's. The design would rather book a detectable, repairable
error than pretend to a prevention it cannot deliver.
